\chapter{Complexity — Counter Prog}
%%% generated by the blueprint pipeline from TCSlib.Complexity.TuringMachine.CounterProg
%%%%%%%%%%%%%%%%%%%%%%%%%%%%%%%%%%%%%%%%%%%%%%%%%%%%%%%%%%%%%%%%%%%%%%%%%%%%%%%
\section{Overview}

This module introduces \emph{counter programs}, the programming model used for the
polynomial-time emitters of [AB09, \S 6.2, Remark 6.7]. A counter program is a goto program
over finitely many registers holding natural numbers. It can increment, decrement (saturating
at $0$) and zero-test a register, print a bit or print a register's value in unary, and read
its binary input once from left to right. The module gives the step semantics of such programs
and compiles every counter program into a binary multi-tape Turing machine with one work tape
per register, holding that register's value in unary. The main result is that one abstract step,
taken while every register is at most $B$, is simulated in at most $2B + 3$ machine steps.

%%%%%%%%%%%%%%%%%%%%%%%%%%%%%%%%%%%%%%%%%%%%%%%%%%%%%%%%%%%%%%%%%%%%%%%%%%%%%%%
\section{Declarations}

\begin{definition}[Instructions of a counter program]
\lean{Complexity.CounterProg.Instr}
\label{Complexity.CounterProg.Instr}
\leanok
Fix a number $R$ of registers, indexed by $\{0, \dots, R-1\}$, and a set $\Lambda$ of labels.
An instruction of a counter program over $R$ registers with labels in $\Lambda$ is one of the
following, where $r$ ranges over registers, $b$ over bits and $\ell', \ell_0, \ell_+, \ell_\bot,
\ell_1$ over labels: $\mathsf{halt}$ (stop); $\mathsf{goto}(\ell')$ (jump to $\ell'$);
$\mathsf{out}(b, \ell')$ (print the bit $b$, then go to $\ell'$); $\mathsf{inc}(r, \ell')$
(increment register $r$); $\mathsf{dec}(r, \ell')$ (decrement register $r$, with $0$ staying
$0$); $\mathsf{jz}(r, \ell_0, \ell_+)$ (go to $\ell_0$ if register $r$ is $0$, and to $\ell_+$
otherwise); $\mathsf{pr}(r, \ell')$ (print $1^v$, where $v$ is the value of register $r$); and
$\mathsf{rd}(\ell_\bot, \ell_0, \ell_1)$ (read the next input symbol: go to $\ell_\bot$ if the
input is exhausted, and otherwise consume the symbol and go to $\ell_0$ or $\ell_1$ according as
it is $0$ or $1$).
\end{definition}

\begin{definition}[Abstract state of a counter program]
\lean{Complexity.CounterProg.St}
\label{Complexity.CounterProg.St}
\leanok
For $R$ registers and a set $\Lambda$ of labels, an abstract state of a counter program consists
of four components: the current label, which is either an element of $\Lambda$ or a special value
marking that the program has halted; a register assignment $\rho : \{0, \dots, R-1\} \to
\mathbb{N}$; the read position $p \in \mathbb{N}$, the number of input symbols consumed so far;
and the output printed so far, a finite bit string.
\end{definition}

\begin{definition}[One step of a counter program]
\lean{Complexity.CounterProg.step}
\label{Complexity.CounterProg.step}
\leanok
\uses{Complexity.CounterProg.Instr, Complexity.CounterProg.St}
Let $P : \Lambda \to \mathrm{Instr}$ be a counter program over $R$ registers (an instruction at
every label) and $x$ an input bit string. One step of $P$ on $x$ leaves a halted state unchanged.
A state with label $\ell$, registers $\rho$, read position $p$ and output $w$ is updated according
to $P(\ell)$, every component not mentioned staying the same:
$\mathsf{halt}$ makes the state halted; $\mathsf{goto}(\ell')$ moves to label $\ell'$;
$\mathsf{out}(b, \ell')$ moves to $\ell'$ and replaces the output by $w b$;
$\mathsf{inc}(r, \ell')$ and $\mathsf{dec}(r, \ell')$ move to $\ell'$ and replace $\rho(r)$ by
$\rho(r) + 1$, respectively by the truncated difference $\max(\rho(r) - 1, 0)$;
$\mathsf{jz}(r, \ell_0, \ell_+)$ moves to $\ell_0$ if $\rho(r) = 0$ and to $\ell_+$ otherwise;
$\mathsf{pr}(r, \ell')$ moves to $\ell'$ and replaces the output by $w\,1^{\rho(r)}$; and
$\mathsf{rd}(\ell_\bot, \ell_0, \ell_1)$ moves to $\ell_\bot$ if $p \ge |x|$, and otherwise moves
to $\ell_0$ or $\ell_1$ according as the symbol $x_p$ (indexing from $0$) is $0$ or $1$, and
increases the read position to $p + 1$.
\end{definition}

\begin{definition}[Running a counter program for $t$ steps]
\lean{Complexity.CounterProg.run}
\label{Complexity.CounterProg.run}
\leanok
\uses{Complexity.CounterProg.Instr, Complexity.CounterProg.St, Complexity.CounterProg.step}
Given a counter program $P$ over $R$ registers with labels in $\Lambda$, an input bit string $x$,
an abstract state $s$ and $t \in \mathbb{N}$, the $t$-step run $\mathrm{run}_{P,x}(s, t)$ is the
state obtained by applying the one-step map of $P$ on $x$ to $s$ exactly $t$ times.
\end{definition}

\begin{definition}[Initial state of a counter program]
\lean{Complexity.CounterProg.init}
\label{Complexity.CounterProg.init}
\leanok
\uses{Complexity.CounterProg.St}
For a label $\ell_0$, the initial abstract state at $\ell_0$ (for $R$ registers) has current
label $\ell_0$, every register equal to $0$, read position $0$, and empty output.
\end{definition}

\begin{theorem}[Zero steps of a counter program]
\lean{Complexity.CounterProg.run_zero}
\label{Complexity.CounterProg.run_zero}
\leanok
\uses{Complexity.CounterProg.St, Complexity.CounterProg.run}
\difficulty{1}
For every counter program $P$ over $R$ registers with labels in $\Lambda$, every input bit string
$x$ and every abstract state $s$, running $P$ on $x$ from $s$ for zero steps returns $s$:
$\mathrm{run}_{P,x}(s, 0) = s$.
\end{theorem}

\begin{theorem}[Runs of a counter program compose]
\lean{Complexity.CounterProg.run_add}
\label{Complexity.CounterProg.run_add}
\leanok
\uses{Complexity.CounterProg.St, Complexity.CounterProg.run}
\difficulty{1}
Let $P$ be a counter program over $R$ registers with labels in $\Lambda$, $x$ an input bit
string and $s$ an abstract state. For all $a, b \in \mathbb{N}$, running $a + b$ steps from $s$
is the same as running $a$ steps and then $b$ further steps:
$\mathrm{run}_{P,x}(s, a + b) = \mathrm{run}_{P,x}(\mathrm{run}_{P,x}(s, a), b)$.
\end{theorem}

\begin{theorem}[Peeling off the first step of a run]
\lean{Complexity.CounterProg.run_succ}
\label{Complexity.CounterProg.run_succ}
\leanok
\uses{Complexity.CounterProg.St, Complexity.CounterProg.run, Complexity.CounterProg.step}
\difficulty{1}
Let $P$ be a counter program over $R$ registers with labels in $\Lambda$, $x$ an input bit
string, $s$ an abstract state and $t \in \mathbb{N}$. Running $t + 1$ steps from $s$ equals
running $t$ steps from the state obtained by one step of $P$ on $x$ from $s$:
$\mathrm{run}_{P,x}(s, t + 1) = \mathrm{run}_{P,x}(\mathrm{step}_{P,x}(s), t)$.
\end{theorem}

\begin{theorem}[A halted state is stationary]
\lean{Complexity.CounterProg.run_of_halted}
\label{Complexity.CounterProg.run_of_halted}
\leanok
\uses{Complexity.CounterProg.St, Complexity.CounterProg.run, Complexity.CounterProg.run_succ, Complexity.CounterProg.step}
\difficulty{4}
Let $P$ be a counter program over $R$ registers with labels in $\Lambda$, $x$ an input bit
string, and $s$ an abstract state whose label component marks it as halted. Then for every
$t \in \mathbb{N}$, running $P$ on $x$ from $s$ for $t$ steps returns $s$ unchanged.
\end{theorem}

\begin{theorem}[One step raises a register by at most one]
\lean{Complexity.CounterProg.step_regs_le}
\label{Complexity.CounterProg.step_regs_le}
\leanok
\uses{Complexity.CounterProg.St, Complexity.CounterProg.step}
\difficulty{4}
Let $P$ be a counter program over $R$ registers with labels in $\Lambda$, $x$ an input bit
string, $s$ an abstract state with registers $\rho$, and $r$ a register. If $\rho'$ denotes the
registers of the state reached from $s$ by one step of $P$ on $x$, then
$\rho'(r) \le \rho(r) + 1$.
\end{theorem}

\begin{theorem}[One step stays within the input]
\lean{Complexity.CounterProg.step_pos_le}
\label{Complexity.CounterProg.step_pos_le}
\leanok
\uses{Complexity.CounterProg.St, Complexity.CounterProg.step}
\difficulty{4}
Let $P$ be a counter program over $R$ registers with labels in $\Lambda$, $x$ an input bit
string, and $s$ an abstract state whose read position $p$ satisfies $p \le |x|$. Then the read
position of the state reached from $s$ by one step of $P$ on $x$ is again at most $|x|$.
\end{theorem}

\begin{theorem}[A $t$-step run raises a register by at most $t$]
\lean{Complexity.CounterProg.run_regs_le}
\label{Complexity.CounterProg.run_regs_le}
\leanok
\uses{Complexity.CounterProg.St, Complexity.CounterProg.run, Complexity.CounterProg.run_succ, Complexity.CounterProg.run_zero, Complexity.CounterProg.step, Complexity.CounterProg.step_regs_le}
\difficulty{4}
Let $P$ be a counter program over $R$ registers with labels in $\Lambda$, $x$ an input bit
string, $s$ an abstract state with registers $\rho$, $r$ a register and $t \in \mathbb{N}$. If
$\rho'$ denotes the registers of the state reached from $s$ after $t$ steps of $P$ on $x$, then
$\rho'(r) \le \rho(r) + t$.
\end{theorem}

\begin{definition}[Control states of the compiled machine]
\lean{Complexity.CounterProg.TSt}
\label{Complexity.CounterProg.TSt}
\leanok
For $R$ registers and a set $\Lambda$ of labels, the control states of the Turing machine
compiled from a counter program are of five kinds, where $r$ is a register and
$\ell, \ell_0, \ell_+$ are labels: $\mathsf{main}(\ell)$, about to execute the instruction at
label $\ell$; $\mathsf{dec}_2(r, \ell)$, the second phase of a decrement of register $r$;
$\mathsf{jz}_2(r, \ell_0, \ell_+)$, the second phase of a zero test of register $r$;
$\mathsf{print}_{\leftarrow}(r, \ell)$, printing register $r$ while walking left over its cells;
and $\mathsf{print}_{\rightarrow}(r, \ell)$, walking back right after such a print. When
$\Lambda$ is finite, so is the set of control states.
\end{definition}

\begin{definition}[Action on a single work tape]
\lean{Complexity.CounterProg.act}
\label{Complexity.CounterProg.act}
\leanok
\uses{Complexity.CounterProg.TSt, Turing.Action}
Consider machines with $R$ work tapes over the bit alphabet whose control states are the compiled
control states for $R$ registers and labels $\Lambda$. Given a work tape index $t$, an input-head
move $i \in \{-1, 0, 1\}$, a write instruction (write nothing, write a blank, or write a bit), a
head move $m \in \{-1, 0, 1\}$, an optional output bit $o$ and an optional next control state
$q$, this is the single-step action that moves the input head by $i$, performs the write
instruction on tape $t$ and moves its head by $m$, leaves every other work tape untouched (no
write, no movement), emits $o$ if present, and passes to $q$ (halting if $q$ is absent).
\end{definition}

\begin{definition}[Action touching no work tape]
\lean{Complexity.CounterProg.ctl}
\label{Complexity.CounterProg.ctl}
\leanok
\uses{Complexity.CounterProg.TSt, Turing.Action}
Consider machines with $R$ work tapes over the bit alphabet whose control states are the compiled
control states for $R$ registers and labels $\Lambda$. Given an input-head move
$i \in \{-1, 0, 1\}$, an optional output bit $o$ and an optional next control state $q$, this is
the single-step action that moves the input head by $i$, neither writes on nor moves any work
tape, emits $o$ if present, and passes to $q$ (halting if $q$ is absent).
\end{definition}

\begin{definition}[Transition table of the compiled machine]
\lean{Complexity.CounterProg.tr}
\label{Complexity.CounterProg.tr}
\leanok
\uses{Complexity.CounterProg.TSt, Complexity.CounterProg.act, Complexity.CounterProg.ctl, Turing.Action}
Let $P$ be a counter program over $R$ registers with labels in $\Lambda$. Its compiled transition
table maps a control state, the symbol under the input head and the symbols under the $R$ work
heads to an action; unless stated otherwise the action writes nothing, moves no head, emits
nothing, and touches only the work tape of the register $r$ involved. In $\mathsf{main}(\ell)$
it acts according to $P(\ell)$: $\mathsf{halt}$ halts; $\mathsf{goto}(\ell')$ passes to
$\mathsf{main}(\ell')$; $\mathsf{out}(b, \ell')$ emits $b$ and passes to $\mathsf{main}(\ell')$;
$\mathsf{inc}(r, \ell')$ writes $1$ on tape $r$, moves its head right and passes to
$\mathsf{main}(\ell')$; $\mathsf{dec}(r, \ell')$, $\mathsf{jz}(r, \ell_0, \ell_+)$ and
$\mathsf{pr}(r, \ell')$ move the head of tape $r$ left and pass to $\mathsf{dec}_2(r, \ell')$,
$\mathsf{jz}_2(r, \ell_0, \ell_+)$ and $\mathsf{print}_{\leftarrow}(r, \ell')$ respectively; and
$\mathsf{rd}(\ell_\bot, \ell_0, \ell_1)$ passes to $\mathsf{main}(\ell_\bot)$ if the input head
reads a blank, and otherwise moves the input head right and passes to $\mathsf{main}(\ell_0)$ or
$\mathsf{main}(\ell_1)$ according as it reads $0$ or $1$. In the auxiliary states, with $c$ the
cell scanned on tape $r$: $\mathsf{dec}_2(r, \ell)$ erases $c$ if it is nonblank and otherwise
moves right, then passes to $\mathsf{main}(\ell)$; $\mathsf{jz}_2(r, \ell_0, \ell_+)$ moves right
and passes to $\mathsf{main}(\ell_+)$ if $c$ is nonblank and to $\mathsf{main}(\ell_0)$
otherwise; $\mathsf{print}_{\leftarrow}(r, \ell)$ emits $1$, moves left and stays put in its
state if $c$ is nonblank, and otherwise moves right and passes to
$\mathsf{print}_{\rightarrow}(r, \ell)$; and $\mathsf{print}_{\rightarrow}(r, \ell)$ moves right
and stays in its state if $c$ is nonblank, and otherwise passes to $\mathsf{main}(\ell)$ without
moving.
\end{definition}

\begin{definition}[Turing machine of a counter program]
\lean{Complexity.CounterProg.toTM}
\label{Complexity.CounterProg.toTM}
\leanok
\uses{Complexity.CounterProg.TSt, Complexity.CounterProg.tr, Turing.FinTM}
Let $\Lambda$ be a finite set of labels, $P$ a counter program over $R$ registers with labels in
$\Lambda$, and $\ell_0 \in \Lambda$ a start label. The machine of $P$ started at $\ell_0$ is the
finite-state multi-tape Turing machine over the bit alphabet with $R$ work tapes (one per
register), whose states are the compiled control states for $R$ and $\Lambda$, whose initial
state is $\mathsf{main}(\ell_0)$, and whose transition function is the compiled transition table
of $P$.
\end{definition}

\begin{definition}[Machine configuration encoding an abstract state]
\lean{Complexity.CounterProg.enc}
\label{Complexity.CounterProg.enc}
\leanok
\uses{Complexity.CounterProg.St, Complexity.CounterProg.TSt, Turing.Cfg, Turing.UnaryTape.ones}
For a natural number $v$ write $1^v$ for the unary tape whose cells $0, \dots, v-1$ hold $1$ and
whose other cells (indexed by $\mathbb{Z}$) are blank. Let $x$ be an input bit string of length
$n$ and $s$ an abstract state with registers $\rho$, read position $p$ and output $w$. Its
encoding is the configuration, on input $x$, of a machine with $R$ work tapes and the compiled
control states, whose state is $\mathsf{main}(\ell)$ if $s$ has label $\ell$ and is halted if $s$
is halted, whose input head is at position $\min(p + 1, n + 1)$, whose work tape $r$ holds
$1^{\rho(r)}$ with its head on cell $\rho(r)$, and whose output tape holds $w$.
\end{definition}

\begin{theorem}[One machine step from a running configuration]
\lean{Complexity.CounterProg.step_some}
\label{Complexity.CounterProg.step_some}
\leanok
\uses{Complexity.CounterProg.TSt, Complexity.CounterProg.toTM, Complexity.CounterProg.tr, Turing.Action.apply, Turing.Cfg, Turing.Cfg.inputSymbol, Turing.MultiTapeTM.step}
\difficulty{1}
Let $\Lambda$ be a finite set of labels, $P$ a counter program over $R$ registers with labels in
$\Lambda$, $\ell_0 \in \Lambda$, and $x$ an input bit string. Let $C$ be a configuration on input
$x$ of the machine of $P$ started at $\ell_0$ which is in some control state $q$ (not halted),
with work tapes $(T_i)$ and work heads $(h_i)$. Then one step of that machine maps $C$ to the
result of applying to $C$ the action that the compiled transition table of $P$ assigns to $q$,
the symbol under the input head of $C$, and the work-tape symbols $(T_i(h_i))_i$.
\end{theorem}

\begin{theorem}[Effect of a single-tape action]
\lean{Complexity.CounterProg.apply_act}
\label{Complexity.CounterProg.apply_act}
\leanok
\uses{Complexity.CounterProg.TSt, Complexity.CounterProg.act, Turing.Action.apply, Turing.Cfg, Turing.UnaryTape.wrT, Turing.moveInputPos}
\difficulty{4}
Let $x$ be an input bit string and consider a configuration on $x$, with $R$ work tapes and the
compiled control states, having an arbitrary state, input-head position $p$, work tapes $(T_i)$,
work heads $(h_i)$ and output $w$. Applying to it the action on the single work tape $t$ with
input-head move $i$, write instruction $\omega$, head move $m$, optional output bit $o$ and
optional next state $q$ yields the configuration with state $q$, input-head position the clamped
move of $p$ by $i$, tape $t$ replaced by $T_t$ with cell $h_t$ overwritten as prescribed by
$\omega$ (and unchanged if $\omega$ writes nothing), head $t$ moved to $h_t + m$, all other
tapes and heads unchanged, and output $w$ followed by $o$ if $o$ is present.
\end{theorem}

\begin{theorem}[Effect of an action touching no work tape]
\lean{Complexity.CounterProg.apply_ctl}
\label{Complexity.CounterProg.apply_ctl}
\leanok
\uses{Complexity.CounterProg.TSt, Complexity.CounterProg.ctl, Turing.Action.apply, Turing.Cfg, Turing.moveInputPos}
\difficulty{3}
Let $x$ be an input bit string and consider a configuration on $x$, with $R$ work tapes and the
compiled control states, having an arbitrary state, input-head position $p$, work tapes $(T_i)$,
work heads $(h_i)$ and output $w$. Applying to it the action that touches no work tape, with
input-head move $i$, optional output bit $o$ and optional next state $q$, yields the
configuration with state $q$, input-head position the clamped move of $p$ by $i$, the same work
tapes $(T_i)$ and heads $(h_i)$, and output $w$ followed by $o$ if $o$ is present.
\end{theorem}

\begin{theorem}[Leftward walk of a register print]
\lean{Complexity.CounterProg.print_walk}
\label{Complexity.CounterProg.print_walk}
\leanok
\uses{Complexity.CounterProg.TSt, Complexity.CounterProg.apply_act, Complexity.CounterProg.step_some, Complexity.CounterProg.toTM, Complexity.CounterProg.tr, Turing.Cfg, Turing.MultiTapeTM.runFrom, Turing.MultiTapeTM.runFrom_succ_eq_step, Turing.MultiTapeTM.runFrom_zero, Turing.UnaryTape.ones, Turing.UnaryTape.ones_lt, Turing.UnaryTape.ones_neg, Turing.UnaryTape.wrT, Turing.moveInputPos_zero}
\difficulty{5}
Let $\Lambda$ be a finite set of labels, $P$ a counter program over $R$ registers with labels in
$\Lambda$, $\ell_0 \in \Lambda$, $x$ an input bit string, $r$ a register, $\ell$ a label, $p$ an
input-head position, $(T_i)$ work tapes and $(h_i)$ work heads, and $m \in \mathbb{N}$ such that
tape $r$ is the unary tape $T_r = 1^m$ (cells $0, \dots, m-1$ hold $1$, all others blank). Then
for every $j \le m$ and every output string $w$, running the machine of $P$ started at $\ell_0$
for $j + 1$ steps from the configuration in state $\mathsf{print}_{\leftarrow}(r, \ell)$ with
input head $p$, tapes $(T_i)$, heads $(h_i)$ except that head $r$ is on cell $j - 1$, and output
$w$, reaches the configuration in state $\mathsf{print}_{\rightarrow}(r, \ell)$ with input head
$p$, tapes $(T_i)$, heads $(h_i)$ except that head $r$ is on cell $0$, and output $w\,1^j$.
\end{theorem}

\begin{theorem}[Rightward walk back after a register print]
\lean{Complexity.CounterProg.back_walk}
\label{Complexity.CounterProg.back_walk}
\leanok
\uses{Complexity.CounterProg.TSt, Complexity.CounterProg.apply_act, Complexity.CounterProg.apply_ctl, Complexity.CounterProg.step_some, Complexity.CounterProg.toTM, Complexity.CounterProg.tr, Turing.Cfg, Turing.MultiTapeTM.runFrom, Turing.MultiTapeTM.runFrom_succ_eq_step, Turing.MultiTapeTM.runFrom_zero, Turing.UnaryTape.ones, Turing.UnaryTape.ones_ge, Turing.UnaryTape.ones_lt, Turing.UnaryTape.wrT, Turing.moveInputPos_zero}
\difficulty{5}
Let $\Lambda$ be a finite set of labels, $P$ a counter program over $R$ registers with labels in
$\Lambda$, $\ell_0 \in \Lambda$, $x$ an input bit string, $r$ a register, $\ell$ a label, $p$ an
input-head position, $(T_i)$ work tapes, $(h_i)$ work heads, $w$ an output string, and
$m \in \mathbb{N}$ such that tape $r$ is the unary tape $T_r = 1^m$ (cells $0, \dots, m-1$ hold
$1$, all others blank). Then for every $d \le m$, running the machine of $P$ started at $\ell_0$
for $d + 1$ steps from the configuration in state $\mathsf{print}_{\rightarrow}(r, \ell)$ with
input head $p$, tapes $(T_i)$, heads $(h_i)$ except that head $r$ is on cell $m - d$, and output
$w$, reaches the configuration in state $\mathsf{main}(\ell)$ with input head $p$, tapes
$(T_i)$, heads $(h_i)$ except that head $r$ is on cell $m$, and output $w$.
\end{theorem}

\begin{theorem}[Updating one unary register tape]
\lean{Complexity.CounterProg.tapes_update}
\label{Complexity.CounterProg.tapes_update}
\leanok
\uses{Turing.UnaryTape.ones}
\difficulty{4}
For $v \in \mathbb{N}$ let $1^v$ be the unary tape whose cells $0, \dots, v-1$ hold $1$ and whose
other cells are blank. Let $\rho : \{0, \dots, R-1\} \to \mathbb{N}$, let $t$ be an index and
$m \in \mathbb{N}$, and let $\rho'$ agree with $\rho$ except that $\rho'(t) = m$. Then replacing
the $t$-th entry of the tape vector $(1^{\rho(r)})_r$ by $1^m$ gives the tape vector
$(1^{\rho'(r)})_r$.
\end{theorem}

\begin{theorem}[Updating one head position]
\lean{Complexity.CounterProg.heads_update}
\label{Complexity.CounterProg.heads_update}
\leanok
\difficulty{4}
Let $\rho : \{0, \dots, R-1\} \to \mathbb{N}$, let $t$ be an index and $m \in \mathbb{N}$, and let
$\rho'$ agree with $\rho$ except that $\rho'(t) = m$. Then replacing the $t$-th entry of the
integer vector $(\rho(r))_r$ by $m$ gives the integer vector $(\rho'(r))_r$.
\end{theorem}

\begin{theorem}[Input head position of an encoded state]
\lean{Complexity.CounterProg.enc_inputPos}
\label{Complexity.CounterProg.enc_inputPos}
\leanok
\uses{Complexity.CounterProg.St, Complexity.CounterProg.enc}
\difficulty{2}
Let $x$ be an input bit string of length $n$ and $s$ an abstract state of a counter program over
$R$ registers whose read position $p$ satisfies $p \le n$. Then in the machine configuration
encoding $s$ on input $x$, the input head is at position $p + 1$.
\end{theorem}

\begin{theorem}[Input symbol read in an encoded state]
\lean{Complexity.CounterProg.enc_inputSymbol}
\label{Complexity.CounterProg.enc_inputSymbol}
\leanok
\uses{Complexity.CounterProg.St, Complexity.CounterProg.enc, Complexity.CounterProg.enc_inputPos, Turing.Cfg.inputSymbol, Turing.FinTM.inputSymbol_at}
\difficulty{1}
Let $x$ be an input bit string of length $n$ and $s$ an abstract state of a counter program over
$R$ registers whose read position $p$ satisfies $p \le n$. Then the symbol under the input head
in the machine configuration encoding $s$ on input $x$ is the bit $x_p$ (indexing from $0$) when
$p < n$, and no symbol (a blank) when $p = n$.
\end{theorem}

\begin{theorem}[One abstract step costs at most $2B+3$ machine steps]
\lean{Complexity.CounterProg.sim_step}
\label{Complexity.CounterProg.sim_step}
\leanok
\uses{Complexity.CounterProg.St, Complexity.CounterProg.TSt, Complexity.CounterProg.act, Complexity.CounterProg.apply_act, Complexity.CounterProg.apply_ctl, Complexity.CounterProg.back_walk, Complexity.CounterProg.enc, Complexity.CounterProg.enc_inputSymbol, Complexity.CounterProg.heads_update, Complexity.CounterProg.print_walk, Complexity.CounterProg.step, Complexity.CounterProg.step_some, Complexity.CounterProg.tapes_update, Complexity.CounterProg.toTM, Complexity.CounterProg.tr, Turing.Action.apply, Turing.Cfg, Turing.MultiTapeTM.runFrom, Turing.MultiTapeTM.runFrom_add, Turing.MultiTapeTM.step, Turing.UnaryTape.ones, Turing.UnaryTape.ones_lt, Turing.UnaryTape.ones_neg, Turing.UnaryTape.update_ones_pred, Turing.UnaryTape.update_ones_succ, Turing.UnaryTape.wrT, Turing.moveInputPos, Turing.moveInputPos_pos_of_ne_right, Turing.moveInputPos_zero}
\difficulty{6}
Let $\Lambda$ be a finite set of labels, $P$ a counter program over $R$ registers with labels in
$\Lambda$, $\ell_0 \in \Lambda$, $x$ an input bit string, and $B \in \mathbb{N}$. Let $s$ be an
abstract state that is running at some label $\ell \in \Lambda$, whose read position is at most
$|x|$, and all of whose registers are at most $B$. Then there is $t \le 2B + 3$ such that running
the Turing machine of $P$ started at $\ell_0$ for $t$ steps from the configuration encoding $s$
on input $x$ yields the configuration encoding the state reached from $s$ by one step of $P$ on
$x$. Here a state with registers $\rho$, read position $p$ and output $w$ is encoded by the
configuration in control state $\mathsf{main}(\ell)$ (or halted if the state is halted) whose
input head is at position $\min(p + 1, |x| + 1)$, whose work tape $r$ holds $1$ exactly on cells
$0, \dots, \rho(r) - 1$ with its head on cell $\rho(r)$, and whose output is $w$.
\end{theorem}
